% GENERATED FILE. Edit canonical lessons, then run npm run generate:books.
% canonical-source-hash: fnv1a64:4f35d847a45b9fa8
% canonical-lessons: FA-C406-qazi, FA-C406-mottaham, FA-C406-shahed, FA-C406-hokm, FA-C406-mohakeme

\chapter{In court}
\label{ch:fa-in-court}
\hlchaptermodality{406}

\begin{chapteropening}
\textbf{By the end of this chapter:} \emph{I can talk about a trial, the judge, the accused and the witnesses, and the verdict.}

Complete the last lesson of chapter 406: I can talk about a trial, the judge, the accused and the witnesses, and the verdict.
\end{chapteropening}

\section[\texorpdfstring{\fa{قاضی} (qâzi)}{قاضی (qâzi)}]{\fa{قاضی} (qâzi) — a judge}

\label{lesson:FA-C406-qazi}



\begin{quote}
Before the new one: say the Persian for a punishment, then the Persian for a permit.
\end{quote}

\subsection*{You'll want to know: \fa{قاضی}}
\textbf{\fa{قاضی}} (\emph{qâzi}) — \textquotedblleft{}a judge\textquotedblright{}.

\begin{quote}
\fa{قاضی} \fa{به} \fa{حرفهای} \fa{هر} \fa{دو} \fa{طرف} \fa{گوش} \fa{داد}.
\end{quote}

\begin{quote}
\emph{qâzi be harfhâ-ye har do taraf gush dâd.}
\end{quote}

\begin{quote}
The judge listened to both sides.
\end{quote}

Plural \emph{qozât}, an Arabic broken plural, or \emph{qâzihâ}.

\begin{cousinweb}
\textbf{\fa{قاضی}} comes from Arabic \emph{qāḍin} \textquotedblleft{}judge\textquotedblright{} (EtymDB); Persian says the Arabic \emph{ḍ} as \emph{z}.
\end{cousinweb}

\begin{grammarlens}[title={What you've built}]
One more word for telling, explaining and describing at B1.
\end{grammarlens}

\subsection*{Your turn}
\begin{itemize}
\raggedright
  \item \emph{Say it:} \emph{qâzi}
  \item \emph{Say it:} \emph{qâzi}, once more
  \item \emph{Say it:} a sentence of your own with \emph{qâzi}
  \item \emph{Recall:} say the Persian for a punishment, then the Persian for a permit, then say \emph{qâzi} again
\end{itemize}

\begin{tcolorbox}[breakable,colback=teal!4,colframe=teal!35!black,title={Before you move on}]
What does \fa{قاضی} mean? (\textquotedblleft{}A judge\textquotedblright{}.) Say it once more.
\end{tcolorbox}
\section[\texorpdfstring{\fa{متهم} (mottaham)}{متهم (mottaham)}]{\fa{متهم} (mottaham) — accused; a defendant}

\label{lesson:FA-C406-mottaham}



\begin{quote}
Before the new one: say the Persian for a permit, then the Persian for a judge.
\end{quote}

\subsection*{You'll want to know: \fa{متهم}}
\textbf{\fa{متهم}} (\emph{mottaham}) — \textquotedblleft{}accused; a defendant\textquotedblright{}.

\begin{quote}
\fa{متهم} \fa{گفت} \fa{که} \fa{آن} \fa{شب} \fa{خانه} \fa{بوده} \fa{است}.
\end{quote}

\begin{quote}
\emph{mottaham goft ke ân shab khâne bude ast.}
\end{quote}

\begin{quote}
The accused said that he had been at home that night.
\end{quote}

Plural \emph{mottahamân}. \emph{Mottaham kardan} — to accuse.

\begin{cousinweb}
\textbf{\fa{متهم}} is usually traced to Arabic \emph{muttaham} \textquotedblleft{}accused\textquotedblright{}, from the same Arabic family as \emph{ittihām} \textquotedblleft{}accusation\textquotedblright{}.
\end{cousinweb}

\begin{grammarlens}[title={What you've built}]
One more word for telling, explaining and describing at B1.
\end{grammarlens}

\subsection*{Your turn}
\begin{itemize}
\raggedright
  \item \emph{Say it:} \emph{mottaham}
  \item \emph{Say it:} \emph{mottaham}, once more
  \item \emph{Say it:} a sentence of your own with \emph{mottaham}
  \item \emph{Recall:} say the Persian for a permit, then the Persian for a judge, then say \emph{mottaham} again
\end{itemize}

\begin{tcolorbox}[breakable,colback=teal!4,colframe=teal!35!black,title={Before you move on}]
What does \fa{متهم} mean? (\textquotedblleft{}Accused; a defendant\textquotedblright{}.) Say it once more.
\end{tcolorbox}
\section[\texorpdfstring{\fa{شاهد} (shâhed)}{شاهد (shâhed)}]{\fa{شاهد} (shâhed) — a witness}

\label{lesson:FA-C406-shahed}



\begin{quote}
Before the new one: say the Persian for a judge, then the Persian for accused.
\end{quote}

\subsection*{You'll want to know: \fa{شاهد}}
\textbf{\fa{شاهد}} (\emph{shâhed}) — \textquotedblleft{}a witness\textquotedblright{}.

\begin{quote}
\fa{دو} \fa{شاهد} \fa{تصادف} \fa{را} \fa{دیدند}.
\end{quote}

\begin{quote}
\emph{do shâhed tasâdof râ didand.}
\end{quote}

\begin{quote}
Two witnesses saw the accident.
\end{quote}

Plural \emph{shâhedân}; the broken plural \emph{shavâhed} means evidence.

\begin{cousinweb}
\textbf{\fa{شاهد}} comes from Arabic \emph{shāhid} \textquotedblleft{}witness\textquotedblright{} (EtymDB), one who was present and saw.
\end{cousinweb}

\begin{grammarlens}[title={What you've built}]
One more word for telling, explaining and describing at B1.
\end{grammarlens}

\subsection*{Your turn}
\begin{itemize}
\raggedright
  \item \emph{Say it:} \emph{shâhed}
  \item \emph{Say it:} \emph{shâhed}, once more
  \item \emph{Say it:} a sentence of your own with \emph{shâhed}
  \item \emph{Recall:} say the Persian for a judge, then the Persian for accused, then say \emph{shâhed} again
\end{itemize}

\begin{tcolorbox}[breakable,colback=teal!4,colframe=teal!35!black,title={Before you move on}]
What does \fa{شاهد} mean? (\textquotedblleft{}A witness\textquotedblright{}.) Say it once more.
\end{tcolorbox}
\section[\texorpdfstring{\fa{حکم} (hokm)}{حکم (hokm)}]{\fa{حکم} (hokm) — a verdict, a ruling}

\label{lesson:FA-C406-hokm}



\begin{quote}
Before the new one: say the Persian for accused, then the Persian for a witness.
\end{quote}

\subsection*{You'll want to know: \fa{حکم}}
\textbf{\fa{حکم}} (\emph{hokm}) — \textquotedblleft{}a verdict, a ruling\textquotedblright{}.

\begin{quote}
\fa{قاضی} \fa{حکم} \fa{را} \fa{خواند}.
\end{quote}

\begin{quote}
\emph{qâzi hokm râ khând.}
\end{quote}

\begin{quote}
The judge read out the verdict.
\end{quote}

Plural \emph{ahkâm}, an Arabic broken plural. Also an order or a decree.

\begin{cousinweb}
\textbf{\fa{حکم}} is usually traced to Arabic \emph{ḥukm} \textquotedblleft{}judgement, ruling\textquotedblright{}, from the same Arabic family as \emph{ḥukūma} \textquotedblleft{}government\textquotedblright{}.
\end{cousinweb}

\begin{grammarlens}[title={What you've built}]
One more word for telling, explaining and describing at B1.
\end{grammarlens}

\subsection*{Your turn}
\begin{itemize}
\raggedright
  \item \emph{Say it:} \emph{hokm}
  \item \emph{Say it:} \emph{hokm}, once more
  \item \emph{Say it:} a sentence of your own with \emph{hokm}
  \item \emph{Recall:} say the Persian for accused, then the Persian for a witness, then say \emph{hokm} again
\end{itemize}

\begin{tcolorbox}[breakable,colback=teal!4,colframe=teal!35!black,title={Before you move on}]
What does \fa{حکم} mean? (\textquotedblleft{}A verdict, a ruling\textquotedblright{}.) Say it once more.
\end{tcolorbox}
\section[\texorpdfstring{\fa{محاکمه} (mohâkeme)}{محاکمه (mohâkeme)}]{\fa{محاکمه} (mohâkeme) — a trial}

\label{lesson:FA-C406-mohakeme}



\begin{quote}
Before the new one: say the Persian for a witness, then the Persian for a verdict.
\end{quote}

\subsection*{You'll want to know: \fa{محاکمه}}
\textbf{\fa{محاکمه}} (\emph{mohâkeme}) — \textquotedblleft{}a trial\textquotedblright{}.

\begin{quote}
\fa{محاکمه} \fa{او} \fa{سه} \fa{ماه} \fa{طول} \fa{کشید}.
\end{quote}

\begin{quote}
\emph{mohâkeme-ye u se mâh tul keshid.}
\end{quote}

\begin{quote}
His trial lasted three months.
\end{quote}

\emph{Mohâkeme kardan} — to put on trial.

\begin{cousinweb}
\textbf{\fa{محاکمه}} is usually traced to Arabic \emph{muḥākama} \textquotedblleft{}trial, bringing to judgement\textquotedblright{}, from the same Arabic family as \emph{ḥukm} \textquotedblleft{}judgement\textquotedblright{}.
\end{cousinweb}

\begin{grammarlens}[title={What you've built}]
One more word for telling, explaining and describing at B1.
\end{grammarlens}

\subsection*{Your turn}
\begin{itemize}
\raggedright
  \item \emph{Say it:} \emph{mohâkeme}
  \item \emph{Say it:} \emph{mohâkeme}, once more
  \item \emph{Say it:} a sentence of your own with \emph{mohâkeme}
  \item \emph{Recall:} say the Persian for a witness, then the Persian for a verdict, then say \emph{mohâkeme} again
\end{itemize}

\begin{tcolorbox}[breakable,colback=teal!4,colframe=teal!35!black,title={Before you move on}]
What does \fa{محاکمه} mean? (\textquotedblleft{}A trial\textquotedblright{}.) Say it once more.
\end{tcolorbox}
